
\chapter{Conclusion}

This thesis aimed to understand how hardware-accelerated ray tracing can be implemented using Vulkan. We explored how acceleration structures, the ray tracing pipeline, and the shader binding table work together to enable hardware-accelerated ray tracing. We also looked at how a mesh can be represented by a bottom-level acceleration structure, and how a node in a scene graph that uses a mesh can be represented as an instance in the top-level acceleration structure. This approach means we do not have to update the entire acceleration structure; instead, we only need to update the meshes that change by rebuilding or refitting their BLAS, and we may need to update the TLAS if the transforms of some nodes have changed.

We also explored how acceleration structures can be updated either by completely rebuilding them or by only refitting the volumes of the BVH to the new positions. In our testing, a combination of both methods—where we refit most of the time and periodically rebuild—proved to be the most efficient in terms of average total rendering time. In many scenarios, this is likely the best approach. However, rendering duration can vary significantly from frame to frame, and a rebuild can cause visible spikes in rendering time. In a real-time application, we would need a way to smooth out these spikes, for example, by rebuilding different meshes across different frames, which is possible thanks to the two-level hierarchy of the acceleration structure. We also have the option to rebuild less dynamic meshes less frequently and more dynamic ones more often, which may further improve performance.
This strategy may not always be ideal. If a large number of rays are cast in the application, it may be better to rebuild the acceleration structures every frame to optimize tracing performance. In these situations, setting the \code{PREFER\_FAST\_TRACE} flag often makes sense.

However, casting only a few rays relative to the number of dynamic objects does not automatically justify refitting every frame, because refitting is not sustainable in many scenes. Refitting only gives good results if the scene does not move too far away from its original state. If the scene keeps changing significantly, rendering time will steadily increase.

Finally, we also saw that not all rays have equal cost. If possible, we can use flags like \code{gl\_RayFlagsSkipClosestHitShaderEXT} and \code{gl\_RayFlagsTerminateOnFirstHitEXT}, which can significantly reduce ray cost. As we observed, our shadow rays were about 37\% faster than regular rays.